\problem{CX01}{Comparisons in Randomized Quicksort}{}
Randomized quicksort sorts $n$ distinct keys. In each nonempty subproblem, it independently chooses a uniform pivot, compares it once with every other key in that subproblem, then recurses on the two sides. Let $C_n$ count all comparisons. Find $\E C_n$ exactly and explain why expected $O(n\log n)$ work does not mean that every run has this order of cost.
\solution Number the keys by increasing rank, $1,\ldots,n$. Let $I_{ij}$ indicate whether ranks $i<j$ are ever compared. Every comparison involves a pivot, which is then removed from future subproblems. A pair can therefore be compared at most once, and
\[
C_n=\sum_{1\le i<j\le n}I_{ij}.
\]
Consider the first pivot selected from ranks $\{i,i+1,\ldots,j\}$. Pivots outside this interval cannot separate $i$ from $j$. If the first pivot inside is an interior rank, it separates the two permanently. If it is either endpoint, the endpoints are compared then. Hence
\[
\Prob(I_{ij}=1)=\frac2{j-i+1}.
\]
Uniformity of the first pivot can be seen through an equivalent implementation: give all keys independent continuous random priorities and always choose the lowest-priority key in a subproblem. This has the same pivot law, and each rank in the interval is equally likely to have its lowest priority.

Linearity of expectation does not require the indicators to be independent. Grouping by $d=j-i$ and writing $H_n=\sum_{k=1}^n1/k$ gives
\[
\begin{aligned}
\E C_n
&=2\sum_{d=1}^{n-1}\frac{n-d}{d+1}\\
&=2(n+1)(H_n-1)-2(n-1)\\
&=\boxed{2(n+1)H_n-4n}.
\end{aligned}
\]
For $n=1$ the formula yields zero; define $C_0=0$ separately. For example, $\E C_4=29/6$.

If every pivot happens to be the smallest key in its subproblem, the cost is $(n-1)+\cdots+1=n(n-1)/2$. This event has positive probability $1/n!$. The expected order thus averages over random runs and is not a worst-case bound for each run.
\insight Ask when a particular pair of keys meets, rather than tracking the whole random recursion tree. Correlated pair indicators can still be added in expectation.

\problem{CX02}{Balancing Shortage and Excess Inventory}{}
Nonnegative demand $D$ has finite mean and distribution function $F$. A real order quantity $q\ge0$ incurs shortage cost $c_u>0$ per unit and excess cost $c_o>0$ per unit. Minimize
\[
L(q)=c_u\E(D-q)_++c_o\E(q-D)_+.
\]
Without assuming a continuous demand distribution, characterize every minimizer. For $\Prob(D=0,1,2)=(1/4,1/2,1/4)$, find all minimizers when $(c_u,c_o)=(3,1)$ and when $(c_u,c_o)=(2,1)$.
\solution For each fixed demand, the loss is a convex, piecewise linear function of $q$, so its expectation is convex. One-sided difference quotients are bounded by $\max(c_u,c_o)$, allowing their limits inside the expectation. Thus
\[
\begin{aligned}
L'_-(q)&=(c_u+c_o)F(q-)-c_u,\\
L'_+(q)&=(c_u+c_o)F(q)-c_u.
\end{aligned}
\]
At $q>0$, convexity makes nonpositive left derivative and nonnegative right derivative necessary and sufficient. Writing $\tau=c_u/(c_u+c_o)$, the full minimizer set is
\[
\boxed{\{q\ge0:F(q-)\le\tau\le F(q)\}}.
\]
At the boundary $q=0$, only the right derivative must be nonnegative. Since nonnegative demand implies $F(0-)=0\le\tau$, the same formula remains valid. Extending the loss formally to $q<0$ makes it strictly decreasing there, so the constraint loses no additional candidate.

A minimizer exists: $\tau\in(0,1)$, and the finite generalized quantile $\inf\{x:F(x)\ge\tau\}$ satisfies the displayed condition. If $F$ is flat at level $\tau$, an entire interval can minimize the loss; uniqueness must not be assumed.

For the discrete example and costs $(3,1)$, $\tau=3/4$. The distribution function jumps to $3/4$ at $1$ and remains there until $2$, yielding all minimizers $[1,2]$, each with expected loss $1$. For costs $(2,1)$, $\tau=2/3$ lies strictly inside the jump at $1$. The unique minimizer is then $q=1$, with minimum loss $3/4$.
\insight Asymmetric absolute loss selects a cost-weighted quantile. The mean is generally not optimal, and atoms or gaps in the distribution determine whether the optimal action is unique.

\problem{CX03}{Maximum Entropy with a Fixed Mean}{}
A nonnegative integer-valued random variable $N$ has mean $\mu>0$. With natural logarithms, define
\[
H(N)=-\sum_{k=0}^\infty p_k\log p_k,\qquad 0\log0=0.
\]
Prove a finite entropy bound and identify the unique maximizing distribution and its entropy. You may use the log-sum inequality for finite distributions, but must handle the infinite support without assuming finite entropy in advance.
\solution Put $r=\mu/(1+\mu)\in(0,1)$ and consider
\[
q_k=(1-r)r^k,\qquad k\ge0.
\]
This geometric distribution has mean $r/(1-r)=\mu$, and its entropy is
\[
H(q)=-\log(1-r)-\mu\log r
=(1+\mu)\log(1+\mu)-\mu\log\mu.
\]

For a rigorous general bound, retain states $0,\ldots,M-1$ separately and merge $N\ge M$ into one tail state of probability $t_M$. The reference distribution assigns probability $r^M$ to this tail. The finite log-sum inequality bounds the grouped entropy by
\[
H_M\le-\sum_{k<M}p_k\log q_k-t_M\log(r^M).
\]
The right-hand side is
\[
-(1-t_M)\log(1-r)
-\log r\left(\sum_{k<M}kp_k+Mt_M\right).
\]
As $M\to\infty$, $t_M\to0$ and the parenthesized quantity is $\E\min(N,M)\uparrow\mu$. Also
$H_M=-\sum_{k<M}p_k\log p_k-t_M\log t_M$ tends to $H(N)$, even if the latter is initially allowed to be infinite, since $-t\log t\to0$. Consequently
\[
H(N)\le(1+\mu)\log(1+\mu)-\mu\log\mu<\infty.
\]

To prove uniqueness, the difference between the grouped cross-entropy bound and $H_M$ is the grouped relative entropy $D_M\ge0$. Its limit is $H(q)-H(N)$. If $p\ne q$, select $j$ with $p_j\ne q_j$. For every $M>j$, merge all groups other than state $j$. The log-sum inequality gives
\[
D_M\ge
p_j\log\frac{p_j}{q_j}
+(1-p_j)\log\frac{1-p_j}{1-q_j}>0.
\]
This positive lower bound does not depend on $M$, excluding equality. Thus the geometric distribution uniquely maximizes entropy.
\insight On an infinite state space, first aggregate the tail and then pass to the limit. This avoids presupposing finite entropy when manipulating relative entropy.
